% =============================================================================
% AI-Assisted GPU Porting of a Large-Scale Weather Simulation Code
% Draft for SC/IPDPS/SE4Science workshop or similar venue
% =============================================================================
\documentclass[conference]{IEEEtran}
% \documentclass[sigconf]{acmart}  % Alternative: ACM format

\usepackage{graphicx}
\usepackage{booktabs}
\usepackage{listings}
\usepackage{xcolor}
\usepackage{hyperref}
\usepackage{amsmath}
\usepackage{algorithm}
\usepackage{algpseudocode}

% Fortran listing style
\lstdefinestyle{fortran}{
  language=[90]Fortran,
  basicstyle=\ttfamily\footnotesize,
  keywordstyle=\color{blue}\bfseries,
  commentstyle=\color{gray}\itshape,
  stringstyle=\color{red},
  numbers=left,
  numberstyle=\tiny\color{gray},
  stepnumber=1,
  numbersep=5pt,
  frame=single,
  breaklines=true,
  captionpos=b,
  morekeywords={acc,kernels,loop,independent,reduction,routine,seq,wait},
}

\lstdefinestyle{bash}{
  language=bash,
  basicstyle=\ttfamily\footnotesize,
  frame=single,
  breaklines=true,
}

\begin{document}

\title{AI-Assisted GPU Porting of a Production Weather Simulation:\\
A Methodology for LLM-Driven OpenACC Migration\\of Large-Scale Fortran/MPI/OpenMP Applications}

\author{
\IEEEauthorblockN{Author A\IEEEauthorrefmark{1},
Author B\IEEEauthorrefmark{1}}
\IEEEauthorblockA{\IEEEauthorrefmark{1}Affiliation\\
Email}
}

\maketitle

% =============================================================================
\begin{abstract}
GPU porting of legacy high-performance computing (HPC) applications remains a
labor-intensive process requiring deep expertise in both the target programming
model and the application domain. We present the first case study of using a
large language model (LLM) code-generation agent to port a production-scale
Fortran/MPI/OpenMP weather simulation code---CReSS (Cloud Resolving Storm
Simulator, approximately 285,000 lines of Fortran across 599 source
files)---to OpenACC for GPU execution. Our methodology decomposes the porting
process into five machine-amenable phases, each designed around the
constraints of current LLM agents: limited persistent memory, bounded context
windows, and non-deterministic output quality. We introduce a verification
pipeline combining per-kernel element-wise validation with a binary search
debugging strategy that isolates faulty GPU kernels in $O(\log N)$ simulation
runs. Applying this methodology, the LLM agent successfully ported 156 GPU
kernels, with four critical bugs identified and fixed through the binary
search process. The resulting GPU implementation achieves a 13.6$\times$
speedup over a 72-thread CPU baseline on an NVIDIA GPU, with numerical
accuracy within floating-point arithmetic ordering differences ($<$1.5\%
variance at 360 time steps). We analyze the types of errors introduced by the
LLM and discuss design principles for human--AI collaborative HPC code
modernization.
\end{abstract}

\begin{IEEEkeywords}
GPU porting, large language models, code generation, OpenACC, Fortran,
weather simulation, HPC code modernization
\end{IEEEkeywords}

% =============================================================================
\section{Introduction}
\label{sec:intro}

The ongoing transition from CPU-centric to GPU-accelerated supercomputing
architectures has created a pressing need to port legacy HPC applications to
GPU programming models. Many production scientific codes, developed over
decades in Fortran with MPI and OpenMP parallelization, require substantial
refactoring to exploit GPU hardware effectively. This porting process is
notoriously labor-intensive: even experienced developers may require months to
years to port a large application, with significant risk of introducing subtle
numerical bugs~\cite{ref:gpu_porting_challenges}.

Large language models (LLMs) have demonstrated remarkable capability in code
generation and translation tasks~\cite{ref:codex, ref:copilot}. However,
their application to HPC code modernization---particularly GPU porting of
large-scale Fortran applications---remains largely unexplored. Existing
studies of LLM-assisted code translation focus on small-scale benchmarks or
single-file programs, leaving open the question of whether LLMs can
effectively contribute to porting production applications with hundreds of
source files and complex inter-module dependencies.

In this paper, we present the first case study of using an LLM code-generation
agent (Claude Code~\cite{ref:claude}) to port a production weather simulation
code from OpenMP to OpenACC. Our target application, CReSS~\cite{ref:cress},
comprises approximately 285,000 lines of Fortran 90 code across 599 source
files, with 387 OpenMP parallel regions. The key challenge lies not in the
individual code transformations---which are well-understood---but in managing
the scale, ensuring correctness, and recovering from the inevitable errors
that LLMs introduce.

Our primary contributions are:

\begin{enumerate}
\item \textbf{A five-phase methodology} for LLM-assisted GPU porting that
  decomposes a monolithic porting task into small, repeatable, verifiable
  units of work, designed around LLM agent constraints
  (Section~\ref{sec:methodology}).

\item \textbf{A verification pipeline} combining per-kernel benchmarking with
  element-wise output comparison and binary search debugging that identifies
  faulty kernels in $O(\log N)$ full-simulation runs
  (Section~\ref{sec:verification}).

\item \textbf{An empirical analysis} of LLM-introduced bugs in GPU porting,
  identifying systematic failure modes including formula truncation,
  copy-paste propagation errors, and conditional branch omission
  (Section~\ref{sec:bugs}).

\item \textbf{Performance and accuracy results} demonstrating 13.6$\times$
  GPU speedup with numerical accuracy within floating-point ordering
  differences (Section~\ref{sec:results}).
\end{enumerate}

% =============================================================================
\section{Background}
\label{sec:background}

\subsection{CReSS Weather Simulation}

CReSS (Cloud Resolving Storm Simulator)~\cite{ref:cress} is a
three-dimensional non-hydrostatic atmospheric model developed at Nagoya
University. It solves the compressible Navier--Stokes equations on a
staggered Arakawa-C grid using a time-splitting method that separates
acoustic (fast) and advective (slow) modes. The code is parallelized with
MPI for domain decomposition and OpenMP for shared-memory threading within
each MPI rank.

Table~\ref{tab:cress_stats} summarizes the code characteristics relevant to
GPU porting. The simulation configuration used for this study employs bulk
cold-rain microphysics (\texttt{cphopt=3}), surface flux parameterization
(\texttt{sfcopt=1}), radiation boundary conditions (\texttt{*bc=7}), and
centered advection (\texttt{advopt=3}) on an $899 \times 899 \times 128$
grid.

\begin{table}[t]
\centering
\caption{CReSS codebase statistics.}
\label{tab:cress_stats}
\begin{tabular}{lr}
\toprule
Metric & Value \\
\midrule
Source files (.f90) & 599 \\
Total lines of code & $\sim$285,000 \\
OpenMP parallel regions & 387 \\
Files containing OpenMP & 338 \\
GPU difficulty: Easy & 236 (61\%) \\
GPU difficulty: Medium & 121 (31\%) \\
GPU difficulty: Hard & 30 (8\%) \\
\bottomrule
\end{tabular}
\end{table}

\subsection{LLM Agent Constraints}

Current LLM code-generation agents operate under several constraints that
fundamentally shape the methodology:

\begin{itemize}
\item \textbf{Context window limitation}: The agent can process a bounded
  amount of text (tokens) per session. For large codebases, this means the
  agent cannot ``see'' the entire application simultaneously. Older context
  is automatically compressed (\emph{auto-compact}) as the conversation
  progresses.

\item \textbf{Session non-persistence}: Each new conversation session starts
  with no memory of previous sessions. Knowledge must be externalized to
  files (progress reports, instruction documents) to persist across sessions.

\item \textbf{Non-deterministic quality}: LLM outputs are probabilistic.
  The same prompt may produce correct code in one invocation and subtly
  buggy code in another. Any methodology must assume that LLM output
  requires verification.
\end{itemize}

These constraints motivate our decomposition into small, independently
verifiable tasks with explicit written specifications---a design pattern we
term \emph{LLM-amenable task decomposition}.

\subsection{Related Work}

LLM-assisted code translation has been studied primarily for high-level
language pairs (e.g., Python to C++, Java to Python)~\cite{ref:llm_translation}.
For HPC-specific tasks, preliminary studies have examined LLM performance on
CUDA kernel generation from natural language descriptions~\cite{ref:llm_cuda}
and OpenMP pragma insertion~\cite{ref:llm_openmp}. However, these studies
typically evaluate on small benchmarks (tens to hundreds of lines) rather
than production applications.

To our knowledge, no prior work has applied LLM agents to the full GPU
porting pipeline of a large-scale Fortran/MPI/OpenMP application, including
analysis, transformation, verification, and debugging at the scale of
hundreds of thousands of lines of code.

% =============================================================================
\section{Methodology}
\label{sec:methodology}

Our methodology decomposes the GPU porting process into five phases, each
producing artifacts that serve as inputs to subsequent phases and as
persistent knowledge across LLM sessions. Figure~\ref{fig:pipeline} provides
an overview.

\begin{figure}[t]
\centering
% TODO: Replace with actual figure
\fbox{\parbox{0.95\columnwidth}{\centering\vspace{2em}
Phase 1: Meta-Info Analysis $\rightarrow$
Phase 2: Runtime Profiling $\rightarrow$\\
Phase 3: CPU Benchmark Extraction $\rightarrow$\\
Phase 4: GPU Benchmark Conversion $\rightarrow$\\
Phase 5: Simulation Integration + Binary Search
\vspace{2em}}}
\caption{Five-phase GPU porting pipeline. Each phase produces persistent
  artifacts (annotated source, profiling data, benchmark programs) that
  externalize knowledge for subsequent LLM sessions.}
\label{fig:pipeline}
\end{figure}

\subsection{Design Principles}

Three principles guide the methodology:

\paragraph{P1: Task atomicity.}
Each task given to the LLM agent must be completable within a single session
and produce a verifiable artifact. A task processes one OpenMP region (one
kernel) at a time, making the workflow inherently parallelizable and
robust to individual failures.

\paragraph{P2: Specification explicitness.}
All instructions are written as structured documents with templates,
checklists, and concrete examples. The LLM agent reads these documents at
the start of each session, effectively receiving a ``fresh briefing'' that
compensates for session non-persistence.

\paragraph{P3: Observable progress.}
Each phase produces human-readable artifacts (annotations, benchmark results,
test logs) that allow the human supervisor to assess progress and quality
without reading the generated code in detail.

\subsection{Phase 1: Meta-Information Analysis}

The LLM agent analyzes each of the 387 OpenMP parallel regions against a
six-item GPU compatibility checklist:

\begin{enumerate}
\item Thread-ID dependencies (\texttt{omp\_get\_thread\_num()})
\item Impure function calls (I/O, global state modification)
\item Shared-variable write conflicts
\item Synchronization primitives (\texttt{critical}, \texttt{atomic}, \texttt{ordered})
\item Reduction operations
\item Other hazards (indirect addressing, complex control flow)
\end{enumerate}

For each region, the agent inserts a structured annotation block directly
into the source file:

\begin{lstlisting}[style=fortran, caption={Meta-information annotation format.}, label={lst:meta}]
!@llm start meta_info
! Location: phvbcs.f90 :: s_phvbcs
! Summary: Scalar boundary phase speed
! GPU diff: Hard
! Findings:
!   - Complex branch structure (4 boundaries)
!   - Conditional dtdvb scaling
! Next: Separate acc kernels per boundary
!@llm end meta_info
\end{lstlisting}

The annotations classify each region as Easy (61\%), Medium (31\%), or Hard
(8\%) for GPU porting. This classification guides subsequent phases and
provides a progress metric.

\subsection{Phase 2: Runtime Profiling}

The LLM agent instruments the simulation to measure the execution time and
call count of each OpenMP region. This profiling data serves two purposes:
(1) prioritizing GPU porting effort toward hot kernels, and (2) providing
the baseline for performance comparison.

\subsection{Phase 3: CPU Benchmark Extraction}

This phase extracts each OpenMP parallel region into a standalone benchmark
program with frozen input/output data. The extraction follows a standardized
template:

\begin{enumerate}
\item \textbf{Variable discovery}: The agent adds \texttt{default(none)} to
  the OpenMP pragma and iteratively compiles until all shared/private
  variables are identified from compiler error messages.

\item \textbf{Data dumping}: Input arrays, scalar parameters, and reference
  output arrays are serialized to binary files at runtime during a short
  simulation run.

\item \textbf{Benchmark creation}: A standalone program reads the dumped data,
  executes the kernel, and validates output against the reference.
\end{enumerate}

Each benchmark directory contains:
\begin{itemize}
\item \texttt{kernel\_benchmark.f90}: Self-contained benchmark program
\item \texttt{data/params.txt}: Scalar parameters
\item \texttt{data/*.bin}: Input and reference output arrays
\item \texttt{benchmark.conf}: Iteration count and validation tolerance
\end{itemize}

This phase produced 413 CPU benchmark directories, covering all OpenMP
regions. The frozen data snapshots decouple kernel development from full
simulation runs, enabling rapid iteration.

\subsection{Phase 4: GPU Benchmark Conversion}

Each CPU benchmark is converted to a GPU (OpenACC) version by the LLM agent
following systematic transformation rules:

\begin{table}[t]
\centering
\caption{OpenMP to OpenACC transformation patterns.}
\label{tab:transform}
\begin{tabular}{p{3.5cm}p{3.5cm}}
\toprule
OpenMP & OpenACC \\
\midrule
\texttt{!\$omp parallel do private(i,j)} &
\texttt{!\$acc kernels} + \texttt{!\$acc loop independent} \\
\texttt{!\$omp do reduction(+:sum)} &
\texttt{!\$acc loop reduction(+:sum)} \\
\texttt{call sub(x)} within parallel region &
\texttt{!\$acc routine seq} on subroutine \\
\bottomrule
\end{tabular}
\end{table}

Data management is handled entirely by NVIDIA Unified Memory
(\texttt{-gpu=managed}), eliminating the need for explicit data directives.
This design choice significantly reduces the complexity of the transformation
task given to the LLM agent, at the cost of potentially suboptimal data
transfer performance.

Each GPU benchmark is validated against the same reference data used by its
CPU counterpart, providing element-wise correctness verification before
integration.

\subsection{Phase 5: Simulation Integration}

The final phase integrates all validated GPU kernels into the main simulation
codebase. Rather than incremental integration (one kernel at a time, $O(N)$
simulation runs), we adopt a \emph{batch integration with binary search
debugging} strategy:

\begin{enumerate}
\item All 156 GPU kernels are integrated simultaneously using preprocessor
  guards:

\begin{lstlisting}[style=fortran, caption={Preprocessor-based kernel control.}, label={lst:guard}]
#if defined(USE_GPU) && !defined(DISABLE_GPU_094)
  ! OpenACC version
  !$acc kernels
  ...
  !$acc end kernels
#else
  ! Original OpenMP (preserved exactly)
  !$omp parallel do ...
  ...
  !$omp end parallel do
#endif
\end{lstlisting}

\item A full simulation test is executed with all kernels enabled.

\item If the simulation fails or produces incorrect results, binary search
  is applied over the kernel space (Section~\ref{sec:verification}).
\end{enumerate}

% =============================================================================
\section{Verification Pipeline}
\label{sec:verification}

Our verification approach addresses a fundamental challenge of LLM-generated
code: the output \emph{appears} correct but may contain subtle semantic
errors invisible to syntactic checking. We employ a three-stage verification
pipeline.

\subsection{Stage 1: Per-Kernel Benchmark Validation}

Each GPU benchmark compares its output element-by-element against the
CPU-generated reference data. This catches most transformation errors at the
individual kernel level. However, certain bugs manifest only in the context
of the full simulation due to:

\begin{itemize}
\item Interaction effects between kernels across time steps
\item Parameter configurations not exercised in the benchmark snapshot
\item Accumulation of small numerical differences over many iterations
\end{itemize}

\subsection{Stage 2: Full Simulation Testing}

The complete simulation is run for a sufficient number of time steps (60--360
steps) and key diagnostic quantities (\texttt{ppmax}, \texttt{ppmin},
\texttt{tkemax}) are compared against the GPU-disabled baseline. A
difference exceeding the expected floating-point ordering variance
($\sim$0.01\% at 60 steps) triggers Stage 3.

\subsection{Stage 3: Binary Search Debugging}

When Stage 2 detects an anomaly, we exploit the per-kernel disable mechanism
to isolate the faulty kernel(s) in $O(\log_2 N)$ simulation runs:

\begin{algorithm}[t]
\caption{Binary search kernel debugging.}
\label{alg:bsearch}
\begin{algorithmic}[1]
\Require Set of kernel IDs $K = \{k_1, k_2, \ldots, k_N\}$, simulation
  test function $\text{Test}(S)$ returning PASS/FAIL for enabled set $S$
\Ensure Faulty kernel ID $k^*$
\State $S \gets K$ \Comment{Start with all kernels enabled}
\If{$\text{Test}(S) = \text{PASS}$}
  \State \Return no fault
\EndIf
\While{$|S| > 1$}
  \State Split $S$ into halves $S_A, S_B$
  \State $R \gets K \setminus S_B$ \Comment{Enable $S_A$, disable $S_B$}
  \If{$\text{Test}(R) = \text{FAIL}$}
    \State $S \gets S_A$ \Comment{Fault is in first half}
  \Else
    \State $S \gets S_B$ \Comment{Fault is in second half}
  \EndIf
\EndWhile
\State \Return $S$ \Comment{Single faulty kernel identified}
\end{algorithmic}
\end{algorithm}

For 156 kernels, this requires at most $\lceil \log_2 156 \rceil = 8$
simulation runs to isolate a single faulty kernel. In practice, 6--7 rounds
sufficed for each of the four bugs discovered (Section~\ref{sec:bugs}).

The binary search is enabled by the preprocessor guard design: each run
requires only recompilation with different \texttt{-DDISABLE\_GPU\_NNN}
flags, not source code changes.

% =============================================================================
\section{LLM-Introduced Bugs: A Taxonomy}
\label{sec:bugs}

Four critical bugs were discovered through the binary search process. All
passed the per-kernel benchmark validation (Stage 1) but caused failures in
full simulation testing (Stage 2). We classify them into three categories.

\subsection{Formula Truncation (Kernel 204: more0q)}

The microphysics source/sink limiter contains complex multi-term formulas for
rain, snow, and graupel precipitation. The LLM omitted terms when
translating the CPU code to GPU:

\begin{itemize}
\item \textbf{Rain sink}: 4 terms missing (\texttt{clrsg}, \texttt{frrg},
  \texttt{shsr}, \texttt{shgr})
\item \textbf{Snow sink}: Sign errors and 3 terms misplaced between
  positive/negative groups
\item \textbf{Graupel sink}: Sign inversion and 3 terms missing
  (\texttt{clri}, \texttt{clir}, \texttt{clsr})
\end{itemize}

This is the most dangerous category: the code compiles, runs, and produces
plausible-looking results. Detection required full simulation testing because
the benchmark data snapshot happened to exercise a parameter regime where the
limiter was not triggered.

\textbf{Root cause}: The formulas span 3--4 lines in the source code,
exceeding the LLM's reliable attention span for faithful term-by-term
copying.

\subsection{Copy-Paste Propagation (Kernels 311, 321)}

Two kernels exhibited errors arising from the LLM incorrectly generalizing
patterns across code blocks:

\paragraph{Kernel 311 (steptund):} The tridiagonal matrix setup requires
three coefficient arrays (\texttt{rr}, \texttt{ss}, \texttt{tt}). The LLM
omitted the \texttt{ss} array entirely in the GPU version, producing only
\texttt{rr} and \texttt{tt}. The uninitialized \texttt{ss} array caused
the Gaussian elimination solver to diverge.

\paragraph{Kernel 321 (swp2nxt):} The time-stepping swap routine has
different semantics for different variable groups: velocity/pressure use
3-way swap (\texttt{past=present; present=future}) while aerosol/tracer/TKE
use 2-way copy (\texttt{past=future}). The LLM applied the 3-way pattern
uniformly to all variables.

\textbf{Root cause}: The LLM recognized a dominant pattern (3-way swap) and
over-generalized it to exceptional cases, a form of \emph{pattern
hallucination}.

\subsection{Conditional Branch Omission (Kernel 094: phvbcs)}

The boundary condition code applies a scaling factor (\texttt{dtdvb}) under
two alternative conditions:

\begin{lstlisting}[style=fortran, caption={Correct conditional logic.}, label={lst:fproc}]
if (fproc(1:3) == 'sml') then
  scpx(j,k,1) = scpx(j,k,1) * dtdvb
else
  if (advopt >= 4) then
    scpx(j,k,1) = scpx(j,k,1) * dtdvb
  end if
end if
\end{lstlisting}

The LLM retained only the \texttt{advopt >= 4} branch, dropping the
\texttt{fproc == 'sml'} condition. Since the simulation uses
\texttt{advopt=3}, the scaling was never applied in the GPU version.

\textbf{Root cause}: The LLM simplified a two-condition branch into one,
possibly interpreting the two branches as redundant since they execute the
same operation.

\subsection{Patterns and Implications}

\begin{table}[t]
\centering
\caption{Summary of LLM-introduced bugs.}
\label{tab:bugs}
\begin{tabular}{llll}
\toprule
ID & Kernel & Category & Detection \\
\midrule
204 & more0q & Formula truncation & 7 rounds \\
094 & phvbcs & Branch omission & 6 rounds \\
311 & steptund & Copy-paste & 7 rounds \\
321 & swp2nxt & Copy-paste & (same run) \\
\bottomrule
\end{tabular}
\end{table}

All four bugs share a common characteristic: they involve \emph{semantic
faithfulness} rather than syntactic correctness. The LLM-generated code is
valid Fortran, compiles without error, and produces outputs within plausible
numerical ranges. This distinguishes LLM bugs from typical human programmer
errors, which more often manifest as compilation failures or crashes.

The bugs also reveal a tension in LLM behavior: the model's tendency to
``simplify'' and ``generalize'' patterns---generally desirable for code
generation---becomes a liability when faithful reproduction of complex
scientific formulas is required.

% =============================================================================
\section{Results}
\label{sec:results}

\subsection{Porting Scale and Effort}

Table~\ref{tab:effort} summarizes the porting effort. The LLM agent
processed 156 GPU kernels across 14 interactive sessions spanning
approximately three weeks. The human effort consisted primarily of designing
the methodology, writing instruction documents, and supervising
binary-search debugging sessions.

\begin{table}[t]
\centering
\caption{Porting effort summary.}
\label{tab:effort}
\begin{tabular}{lr}
\toprule
Metric & Value \\
\midrule
GPU kernels ported & 156 \\
CPU benchmarks created & 413 \\
GPU benchmarks created & 172 \\
LLM sessions & 14 \\
Calendar time & $\sim$3 weeks \\
Git commits & 32 \\
Bugs found by binary search & 4 \\
Binary search rounds (total) & $\sim$20 \\
\bottomrule
\end{tabular}
\end{table}

\subsection{Numerical Accuracy}

Table~\ref{tab:accuracy} compares the GPU simulation output against two
baselines: (1) the CPU-only build without \texttt{-acc} flags
(\emph{noacc}), and (2) the GPU build with all GPU kernels disabled
(\emph{gpudisabled}). The gpudisabled baseline isolates the effect of
compiler flags from GPU kernel execution.

\begin{table}[t]
\centering
\caption{Numerical accuracy comparison at 60 and 360 time steps.
  Differences are relative to the gpudisabled baseline.}
\label{tab:accuracy}
\begin{tabular}{lrr}
\toprule
Diagnostic & 60 steps & 360 steps \\
\midrule
\texttt{ppmax} & $+0.000$\% & $+0.000$\% \\
\texttt{ppmin} & $+0.000$\% & --- \\
\texttt{tkemax} & $+0.001$\% & $-1.464$\% \\
\bottomrule
\end{tabular}
\end{table}

At 60 time steps, the GPU-enabled simulation is virtually identical to the
GPU-disabled baseline. At 360 steps, \texttt{tkemax} shows a 1.464\%
difference attributable to floating-point arithmetic ordering differences
between CPU (sequential reductions) and GPU (parallel reductions). This
level of variance is consistent with known GPU porting behavior and does not
indicate algorithmic error.

Separately, we confirmed that the compiler flags themselves
(\texttt{-acc -gpu=managed}) introduce negligible numerical differences:
\texttt{ppmax} differs by $\sim 3.5 \times 10^{-7}$ between the noacc and
gpudisabled builds.

\subsection{GPU Performance}

% TODO: Add detailed performance data when available
% Preliminary results show 13.6x overall speedup
The GPU implementation achieves a preliminary overall speedup of
approximately 13.6$\times$ compared to the 72-thread OpenMP CPU baseline
for a 360-step simulation. Individual kernel speedups vary widely, from
$\sim$2$\times$ for memory-bandwidth-limited stencil operations to over
100$\times$ for embarrassingly parallel operations
(Table~\ref{tab:performance}).

\begin{table}[t]
\centering
\caption{Representative kernel speedups (GPU vs.\ 72-thread CPU).}
\label{tab:performance}
\begin{tabular}{llr}
\toprule
Kernel & Description & Speedup \\
\midrule
outmxn & Min/max diagnostics & 187$\times$ \\
adjstni & Ice concentration adjust & 22$\times$ \\
setblk & Block initialization & 22$\times$ \\
swp2nxt & Time level swap & 19$\times$ \\
\bottomrule
\end{tabular}
\end{table}

% =============================================================================
\section{Cost Model and Stability}
\label{sec:cost}

A critical but often overlooked aspect of LLM-assisted software engineering
is the \emph{cost of failure}. Unlike interactive development where a human
programmer can immediately recognize and fix errors, LLM-generated code
failures trigger expensive re-execution cycles involving GPU compute time,
LLM token consumption, and human supervisory attention. This section
analyzes the stability characteristics of each phase and the design
decisions that mitigate rework costs.

\subsection{The Rework Cost Model}

Each failure in the pipeline incurs three costs:

\begin{itemize}
\item \textbf{Compute cost}: Full simulation runs require 2--4 hours of GPU
  time; even benchmark runs consume significant resources when repeated
  across hundreds of kernels.
\item \textbf{Token cost}: LLM sessions consume tokens proportional to context
  size. Debugging sessions that revisit previously generated code effectively
  ``pay twice'' for the same transformation.
\item \textbf{Attention cost}: Human supervision time is the scarcest
  resource. Each failure requires the human to diagnose whether the issue is
  in the LLM output, the specification, or the verification infrastructure.
\end{itemize}

Table~\ref{tab:rework} summarizes the rework observed across phases.

\begin{table}[t]
\centering
\caption{Rework statistics by phase. Phase 3 required the most iterations
  due to output format inconsistency.}
\label{tab:rework}
\begin{tabular}{lrrr}
\toprule
Phase & Fix commits & Re-runs & Calendar time \\
\midrule
1: Meta-info analysis & 2 & 0 & 2 days \\
2: Runtime profiling & 1 & 1 & 1 day \\
3: CPU benchmark extraction & 11+ & 15+ & $\sim$6 weeks \\
4: GPU benchmark conversion & 3 & 5 & $\sim$2 weeks \\
5: Integration + debug & 4 & $\sim$20 & $\sim$2 weeks \\
\bottomrule
\end{tabular}
\end{table}

\subsection{Phase 3: The Bottleneck}

CPU benchmark extraction (Phase~3) was by far the most problematic phase,
consuming approximately half of the total project calendar time. The root
cause was \emph{output format drift}: while the LLM consistently performed
the requested operation (dumping data to files), the specific format varied
across invocations in ways that broke downstream consumption.

Specific failure modes included:

\paragraph{Data serialization inconsistency.}
The LLM was instructed to dump scalar parameters to a text file
(\texttt{params.txt}). Across different sessions, it produced three
incompatible formats: (1)~positional values only, (2)~\texttt{key = value}
pairs, and (3)~Fortran namelist format. Since the benchmark reader was
generated in a separate session, format mismatches caused silent data
corruption or runtime errors.

\paragraph{Conditional allocation blindness.}
The simulation uses compile-time flags (\texttt{cphopt}, \texttt{sfcopt},
\texttt{savmem}) to conditionally allocate arrays. The LLM-generated dump
code attempted to serialize arrays that were not allocated under the current
configuration, causing segmentation faults. Six separate fix commits over
four days were required to address cascading instances of this issue.

\paragraph{Type confusion.}
The LLM called scalar dump routines (\texttt{dump\_scalar\_r()}) on array
variables, causing silent data truncation. A single fix commit modified
140~source files with 1,434 insertions and 540 deletions to correct
variable type handling across all kernels.

\paragraph{Endianness mismatch.}
The main simulation was compiled with \texttt{-Mbyteswapio} (big-endian I/O),
but the LLM-generated benchmarks did not include the corresponding flag,
producing unreadable binary data.

\subsection{Mitigation: Template Rigidity vs.\ LLM Creativity}

The Phase~3 difficulties reveal a fundamental tension in LLM-assisted code
generation: the model's tendency to produce \emph{creative variations} of
the requested output, while the downstream pipeline requires \emph{exact
format compliance}. We addressed this through progressively stricter
templates:

\begin{enumerate}
\item \textbf{Shared module}: A standardized \texttt{dump\_kernel\_data.f90}
  module with typed dump routines (\texttt{dump\_scalar\_r},
  \texttt{dump\_array\_3d}, etc.) eliminated format variation in data
  serialization.

\item \textbf{Rigid directory structure}: Each benchmark directory must
  contain exactly \texttt{kernel\_benchmark.f90}, \texttt{Makefile},
  \texttt{benchmark.conf}, and \texttt{data/params.txt}, with prescribed
  formats for each.

\item \textbf{Makefile inclusion}: GPU benchmarks use
  \texttt{include ../Makefile.common} to eliminate compiler flag variation.

\item \textbf{Explicit variable discovery}: The \texttt{default(none)}
  compilation technique forces exhaustive variable enumeration, preventing
  omissions.
\end{enumerate}

These mechanisms transformed Phase~3 from a high-variance process to a
mechanical one. The key insight is that \textbf{the template must constrain
the LLM's output space to a single valid format}, leaving no room for
``equivalent'' alternatives that break downstream compatibility.

\subsection{Stability Across Phases}

Figure~\ref{fig:stability} illustrates the relationship between template
rigidity and phase stability. Phases with well-defined templates (1, 4, 5)
required minimal rework, while Phase~3, where the output format was
initially under-specified, required extensive iteration.

\begin{figure}[t]
\centering
% TODO: Replace with actual figure
\fbox{\parbox{0.95\columnwidth}{\centering\vspace{2em}
Scatter plot: X = template rigidity score (1--5),\\
Y = number of fix commits per phase.\\
Phase 3 is a clear outlier at (low rigidity, high fixes).\\
After template hardening, subsequent kernels required no fixes.
\vspace{2em}}}
\caption{Template rigidity vs.\ rework effort. Phase~3 initially had
  low template rigidity, resulting in high rework. After introducing
  the shared dump module and rigid directory structure, stability improved
  dramatically.}
\label{fig:stability}
\end{figure}

% =============================================================================
\section{Discussion}
\label{sec:discussion}

\subsection{Prompt Engineering for HPC Code Transformation}

The success of our methodology relies heavily on \emph{structured prompt
design} that accommodates LLM limitations:

\paragraph{Task atomicity and loop processing.}
Each phase is designed so that the LLM processes one kernel at a time in a
loop: read specification $\rightarrow$ read source $\rightarrow$ generate
output $\rightarrow$ validate $\rightarrow$ next kernel. This pattern ensures
that each task fits within the context window and that progress is preserved
even if the session terminates unexpectedly.

\paragraph{Template-driven output.}
By providing rigid output templates (annotation format, benchmark directory
structure, preprocessor guard pattern), we constrain the LLM's output space
and reduce the opportunity for creative but incorrect variations.

\paragraph{Externalized memory.}
Progress reports, instruction documents, and the \texttt{CLAUDE.md} memory
file serve as the LLM's ``persistent memory,'' enabling coherent work across
sessions despite the lack of built-in session persistence. This creates a
form of \emph{software-defined memory} that compensates for the LLM's
architectural limitations.

\subsection{Human--AI Role Separation}

Our experience suggests a natural role separation:

\begin{itemize}
\item \textbf{Human}: Methodology design, specification writing, simulation
  configuration, binary search supervision, final validation judgment
\item \textbf{LLM}: Code analysis, benchmark extraction, OpenMP-to-OpenACC
  transformation, bug fixing under direction
\end{itemize}

The human provides \emph{architectural} decisions (what to do), while the
LLM provides \emph{operational} execution (how to do it at scale). This
mirrors traditional software engineering roles but at a different
granularity.

\subsection{Benchmark Infrastructure as a Tuning Pipeline}

A significant secondary benefit of the per-kernel benchmark framework is
that it provides a \emph{ready-made performance tuning infrastructure}. Each
GPU benchmark is a self-contained program that can be independently compiled,
profiled, and optimized without running the full simulation. This enables:

\begin{itemize}
\item \textbf{Isolated kernel tuning}: OpenACC directives (e.g.,
  \texttt{collapse}, \texttt{tile}, \texttt{gang/worker/vector} clauses)
  can be tested on individual kernels with immediate performance feedback.
\item \textbf{Rapid iteration}: Benchmark compilation takes seconds vs.\
  minutes for the full simulation, enabling orders-of-magnitude faster
  optimization cycles.
\item \textbf{Correctness-preserving optimization}: After tuning, the
  benchmark validates output against the frozen reference data, ensuring
  that optimizations do not introduce numerical errors.
\item \textbf{Re-integration}: Once a kernel is tuned and validated in
  isolation, the optimized code is copied back into the simulation source
  within the existing preprocessor guard structure.
\end{itemize}

This design creates a virtuous cycle: the verification infrastructure built
for correctness (Phase~3--4) becomes the optimization infrastructure for
performance (post-Phase~5). The LLM agent can also be employed for the
tuning phase, applying the same template-driven workflow to generate
optimization variants.

\subsection{Comparison with Manual Porting}

While we do not have a controlled comparison with purely manual porting, we
note that the three-week calendar time for 156 kernels across 285K lines of
code is competitive with published GPU porting efforts for comparable
applications, which typically report months of developer effort. The
methodology's main advantage is not necessarily speed but
\emph{systematization}: the structured pipeline reduces the risk of ad-hoc
errors and ensures comprehensive coverage.

\subsection{Limitations and Threats to Validity}

\paragraph{Single LLM evaluation.}
This study used a single LLM (Claude Code with Opus model). The
methodology's effectiveness may vary across different models. LLMs differ in
their Fortran proficiency, tendency toward output format drift, and ability
to faithfully reproduce complex formulas. Cross-model evaluation using
alternative agents such as OpenAI Codex/GPT-4, open-source models (Code
Llama, DeepSeek-Coder), or specialized HPC code models would strengthen the
generalizability claims. Importantly, our methodology is \emph{model-agnostic}
by design: the five-phase pipeline, template structure, and binary search
verification are independent of the underlying LLM. The per-kernel benchmark
infrastructure provides a natural framework for comparative evaluation---the
same benchmark suite could validate GPU kernels generated by different models.

\paragraph{Data management.}
Our use of unified memory (\texttt{-gpu=managed}) sacrifices potential
performance gains from explicit data management. Future work could apply the
same methodology to optimize data transfers, using the benchmark
infrastructure for isolated profiling of data movement patterns.

\paragraph{Benchmark coverage.}
The binary data snapshot captures only one parameter configuration. Bugs that
manifest under different configurations may escape benchmark validation, as
observed with Kernel~204 (more0q), where the microphysics limiter was not
triggered in the benchmark snapshot. Multi-configuration benchmark generation
could mitigate this at the cost of increased data storage and setup
complexity.

\paragraph{Single application.}
Our methodology was developed and tested on a single application.
Generalizability to other HPC codes---particularly those with different
parallelization patterns (e.g., task-based parallelism, non-structured
grids)---requires further validation.

\paragraph{Performance optimization.}
The current work focuses on correctness. Systematic GPU performance
optimization (kernel fusion, data locality, explicit data management)
remains as future work, for which the per-kernel benchmark infrastructure
provides a ready-made foundation (Section~\ref{sec:discussion}).

% =============================================================================
\section{Lessons Learned}
\label{sec:lessons}

\begin{enumerate}
\item \textbf{Verification infrastructure is the most valuable investment.}
  The per-kernel benchmark framework and binary search debugging capability
  consumed significant setup effort but were essential for achieving correct
  results. Without them, the four subtle LLM-introduced bugs would have been
  extremely difficult to diagnose.

\item \textbf{LLMs struggle with faithful formula reproduction.} Multi-term
  scientific formulas should be treated as high-risk transformation targets.
  Additional verification mechanisms (e.g., term-counting assertions, symbolic
  comparison) could mitigate this.

\item \textbf{The ``simplification instinct'' of LLMs is dangerous for
  scientific code.} LLMs tend to merge similar-looking code branches and
  normalize patterns. This is often correct for business logic but can
  silently change the semantics of physics-based computations.

\item \textbf{Preprocessor-based kernel control enables rapid debugging.}
  The ability to enable/disable individual GPU kernels via compile flags,
  without modifying source code, was crucial for the binary search strategy
  and should be adopted as a standard practice for GPU porting projects.

\item \textbf{Session boundaries require careful knowledge management.} The
  most error-prone moments in the project were session transitions, where
  context was lost. Investing in comprehensive progress reports and
  instruction documents paid significant dividends.

\item \textbf{Output format stability is the single largest risk factor.}
  Phase~3 consumed half the project time due to format drift in
  LLM-generated data serialization code. The lesson is stark: any interface
  between LLM-generated components must be constrained by a shared,
  machine-enforced format specification (e.g., a shared module), not by
  natural-language instructions alone. Telling the LLM ``dump data in binary
  format'' is insufficient; providing a \texttt{dump\_array\_3d()} routine
  that the LLM must call is robust.

\item \textbf{Cost-aware phase ordering matters.} Phases should be ordered
  so that the cheapest verification comes first. Per-kernel benchmarks
  (seconds to run) catch most bugs before full simulation testing (hours to
  run), and binary search (logarithmic runs) avoids exhaustive testing.
  This ``verification funnel'' minimizes total compute cost.
\end{enumerate}

% =============================================================================
\section{Conclusion}
\label{sec:conclusion}

We have presented the first case study of LLM-assisted GPU porting for a
production-scale Fortran/MPI/OpenMP weather simulation code. Our five-phase
methodology, designed around the constraints of current LLM agents,
successfully ported 156 GPU kernels in a 285,000-line codebase, achieving
13.6$\times$ GPU speedup with numerical accuracy within floating-point
ordering differences.

The key insight is that LLMs are effective at the \emph{mechanical} aspects
of GPU porting (syntactic transformation, boilerplate generation, pattern
application) but unreliable for \emph{semantic} faithfulness (formula
reproduction, conditional logic preservation). A verification pipeline that
assumes LLM output is ``probably correct but possibly wrong'' is essential.

Our binary search debugging strategy, enabled by preprocessor-based kernel
control, provides a practical and efficient approach to identifying the
inevitable errors. The four bugs discovered---formula truncation, copy-paste
propagation, and conditional branch omission---represent systematic failure
modes that tool developers can target for improvement.

We believe this methodology is applicable to other large-scale HPC code
modernization efforts and represents a practical model for human--AI
collaboration in scientific software engineering.

% =============================================================================
% References (placeholder - replace with actual bibtex entries)
\begin{thebibliography}{99}

\bibitem{ref:cress}
K.~Tsuboki and A.~Sakakibara, ``Large-scale parallel computing of cloud
resolving storm simulator,'' in \emph{High Performance Computing}, Springer,
2002, pp.~243--259.

\bibitem{ref:claude}
Anthropic, ``Claude Code: An agentic coding tool,'' 2025. [Online].
Available: https://docs.anthropic.com/en/docs/claude-code

\bibitem{ref:gpu_porting_challenges}
S.~Memeti et al., ``Benchmarking OpenCL, OpenACC, OpenMP, and CUDA:
Programming productivity, performance, and energy consumption,'' in
\emph{Proc.\ Workshop on Adaptive Resource Management and Scheduling for
Cloud Computing (ARMS-CC)}, 2017.

\bibitem{ref:codex}
M.~Chen et al., ``Evaluating large language models trained on code,''
\emph{arXiv preprint arXiv:2107.03374}, 2021.

\bibitem{ref:copilot}
A.~Ziegler et al., ``Productivity assessment of neural code completion,''
in \emph{Proc.\ 6th ACM SIGPLAN Int.\ Symp.\ on Machine Programming}, 2022.

\bibitem{ref:llm_translation}
B.~Roziere et al., ``Code Llama: Open foundation models for code,''
\emph{arXiv preprint arXiv:2308.12950}, 2023.

\bibitem{ref:llm_cuda}
N.~Nichols et al., ``Can large language models write parallel code?,'' in
\emph{Proc.\ 33rd Int.\ Symp.\ on High-Performance Parallel and Distributed
Computing (HPDC)}, 2024.

\bibitem{ref:llm_openmp}
T.~Kadosh et al., ``Domain-specific LLMs for parallel program generation,''
\emph{arXiv preprint arXiv:2401.10062}, 2024.

\end{thebibliography}

\end{document}
